\chapter{Teaching Material and the Instruction Corpus}
\label{ch:teaching}

\section{A Third Party in the Division of Labour}

The paper has so far described two parties. There is a third, and it is easy to
miss because it does not run: everything an agent reads before it authors anything.
That is two bodies of material, the worked examples and the instructions, and this
chapter takes them in that order.

This matters more here than it would in most systems, because an agent authoring
blind cannot learn a primitive by trying it and looking. It learns by reading an
example, and what it reads shapes what it builds. A teaching corpus that
demonstrates one instrument per map teaches an agent to reach for one instrument
per map. The repository has learned this the expensive way and has changed its
teaching format accordingly, which makes the change itself worth documenting.

\section{Two Formats, and What Each Holds Constant}

The repository carries two bodies of worked examples.

The older is a \textbf{showcase} of twenty-three boards, one technique each. Its
method is to hold everything else fixed: every board forks one control board, a
plain hundred-by-hundred destroy board that scores zero with no violation and no
complaint, and changes \emph{only} what its technique needs. The diff is therefore
the lesson, with nothing else in it.

The newer is a set of six \textbf{technique cards}. Its method is more radical: it
deletes the board. A technique card has no teams, no spawns, no objectives, no
gamemode and no plan at all. Its whole document is a setup, its themes, a map theme
and its layers. There is nothing to diff against because there is nothing else
present.

The size difference follows from that and is not incidental. The six cards come to
392\,KB. The twenty-three showcase boards come to 12\,MB, because each carries a
built world. The showcase's own README promises examples short enough to read in
full and small enough to paste; the cards are what actually delivers it, with four
complete outlines stated in 4{,}395 bytes.

\section{Five Differences That Change What Is Learned}

\subsection{Removing the Symmetry Confounder}

Every technique card states \id{mirror\_mode: "none"} and every group in it states
that it does not mirror. The showcase cannot: it is made of real boards, so it must
spend a section of its README on what the symmetry fan decides, and it records
three boards whose pieces landed in each other because of it, at the cost of a
rebuild.

Mirroring is a real subject and an agent must learn it. The point is that it
becomes a subject one can learn \emph{separately}, rather than a tax on learning
anything else. The relief card still states the fact, that any symmetry mode folds
the relief across the centre and a group's own setting does not change it. It
simply does not make the reader pay for it while learning about relief.

\subsection{Comparison Inside One Render}

This is the largest pedagogical change. The showcase is one technique per map, so a
reader comparing two treatments reads two folders and holds the difference in their
head. A card puts the variants side by side in one world: five flights up the same
eight blocks; the same island flat and at four heights; three relief plots
sixty-four blocks apart; four outlines in a row.

The controlled comparison is then \emph{in the picture}, and the relief card takes
it further by making it quantitative. Its two plots differ only in the heights their
shapes are drawn at, and the committed check reads: every one of one plot's 2{,}304
columns has the same top as the column sixty-four blocks east of it, and none
differ. That is a proof that a solved shape's drawn height decides nothing, rather
than an assertion of it.

\subsection{The One Mistake as a Standing Section}

Five of the six cards carry a heading titled, literally, \emph{The one mistake}. It
names the failure, its rule id, and the threshold at which it fires. The set is a
good summary of where a terrain author actually goes wrong:

\begin{table}[htbp]
\caption{The mistake each card guards against, with the threshold where each one
begins to bite.}
\label{tab:one-mistake}
\small
\begin{tabularx}{\textwidth}{L{3.5cm} L{1.6cm} Y}
\toprule
\textbf{Card} & \textbf{Rule} & \textbf{The mistake} \\
\midrule
combined shapes & \id{SK9} & a lower shape drawn inside a higher one on the same
  layer is not in the world \\
polylines & \id{SK25} & a band that turns tighter than it is wide laps itself.
  Four points across 26 blocks at radius 3 lap in every bend; the same four spread
  over 40 clear \\
curved outlines & --- & a handle pulling further from its edge than along it draws
  a loop rather than a bulge \\
ramp and stair & --- & whole-number anchors 8 and 1 over 16 blocks build treads 1,
  2 and 3 deep. A flight comes out even only with anchors half a riser past its
  first and last tread \\
relief on shapes & --- & a made shape with no \id{relief\_scope} joins the solve.
  The same yard reads $y11$ instead of $y17$ \\
\bottomrule
\end{tabularx}
\end{table}

The showcase scatters the equivalent knowledge through prose and collects it in a
thirty-row index at the end. Making it structural and local is what lets an agent
find it at the moment it is about to make the mistake.

\subsection{Verification Committed as Text}

The repository's own rule is to read the text before the pictures, and the cards
implement it literally. Every card ends with a section naming the committed read
that proves each of its claims: a census, a column, a section, a heightmap, a set
of transects, a walk, a slope grid, a comparison, a relief read.

Those are the API's actual answers, verbatim, re-checkable, and diffable across a
studio upgrade. The showcase's verification is a folder of pictures plus a table of
numbers typed into its README, which cannot be re-run.

The strongest instance is the relief card's comparison file, which prints the same
rows three times: as drawn, with \id{hold}, and with \id{exclude}. The
counterfactual is committed alongside the fact, which is a standard of evidence
most documentation does not attempt.

\subsection{Teaching the Revision Surface}

The showcase is document-in, world-out: one call drives a whole spec. That teaches
an agent to author, and teaches it nothing about \emph{revising}.

One technique card is the only place in either body that demonstrates the
incremental editing surface, with every request and its response committed: moving
a vertex, inserting one, deleting one, and setting a shape's anchor heights. It
also captures the arithmetic that bites, which is that an insert lands one past the
index it names and every later index moves up by one. An agent replaying a list of
edits computed against the original ring will scramble the outline, and this is the
only example anywhere that would have warned it.

\section{What the Older Format Still Holds}

The honest framing is narrower than \emph{the cards replace the showcase}. The
cards cover geometry and paint, on ground with nothing on it. The showcase covers
the entire non-terrain half of the system: objectives and their gate rules,
plan-tier authoring, walls, dressing, room styles, coverage, traversability,
symmetry and the export gate.

So the cards supersede the showcase for the terrain-authoring primitives, and the
showcase remains the record for everything board-shaped. That is already how they
are used in practice: one orchestrated run, briefing twelve authoring agents,
records giving them the cards in preference to the showcase, alongside a ground
family chosen to stop twelve boards coming out as twelve greys.

\section{The Instruction Corpus}

The worked examples are one half of what an agent reads. The other half is a body
of instructions, and it is worth describing because its shape has been argued over
as carefully as the cards' and for the same reason.

There are seven documents, coming to about 3{,}800 lines. Four are briefs, naming
what a board should be: a general authoring brief, a revamp brief for changing an
existing board, a brief on sculpting with layers, and one on adapting a composed
board. One is a notes file recording what the API still cannot be asked. Two are
skills, loaded by name rather than read as prose, and they are the ones that carry
the procedure.

\subsection{Two Skills, Loaded in Order}

The first skill is a warm-up, loaded once at the start of a run and before
anything else. Its subject is what reading costs. The documents an agent is
pointed at come to roughly 82{,}000 tokens, which is 41\% of a two-hundred-thousand
token window spent before a single map is opened, and the skill's first table is a
budget: what each read costs and when it may be taken. It carries a never-open
list. Two stored layouts exceed a whole context window at 397{,}000 and 367{,}000
tokens, and the skill states plainly that opening either one ends the run.

The rule that follows from the budget is that a layout is queried rather than read,
and that a one-line answer is a search rather than a file read. A slope grid runs
to ten thousand tokens and its verdict is one line.

The second skill is the authoring procedure, loaded whenever a board is built,
changed, diagnosed or read. It opens with a lookup table of about thirty rows, and
the table has three columns: the question, the read that answers it, and
\emph{what not to use instead}. The third column is the unusual one and is the
reason the table works.

\subsection{The Third Column}

A row of that table reads: \emph{how steep is the ground, and where} is answered by
the incline report, and \emph{not} by the slope grid, which answers whether a step
can be walked and not what angle the ground lies at. Another reads: \emph{is the
shape I authored in the world at all} is answered by a column read, and not by the
store's 200 or the export gate opening, neither of which looks.

Each of those negatives is a specific mistake that happened. The corpus of
seventy-one run reports behind the skill records the most-repeated sentence in it
as a variant of \emph{the instrument I should have used exists and I did not know
about it}, and the second most common as using a read that answers a neighbouring
question. A table that named only the right read would not prevent the second.

\subsection{The Two Moments}

The skill names two places to stop, and that is all it names. The first is after
the plan is written and before the first world is built: render the plan as a grid
and read it, because a plan is a list of rectangles and most of what goes wrong
with one is a relation between two of them, which no render of a built world can
show, because by then they are terrain.

The second is after every build and before any picture is opened: three text reads,
named individually, whose numbers are to be said out loud. The ordering is stated
as a rule rather than a preference, in a sentence the repository repeats
everywhere: a picture answers whether a thing came out, and a number answers
whether it is right. The example attached to it is a one-block bump under a rail,
which is one shade in a heightmap and nothing at all in an isometric, and which
shipped in five consecutive builds of one board because nobody opened the file that
named it with coordinates.

\subsection{One Paragraph, One Claim}

The instruction corpus has its own writing rule, and it is enforced by a script
rather than by review.

An instruction is not read the way an essay is read. It is scanned for the sentence
that applies right now, so a paragraph carrying six claims hides five of them. The
failure has one shape: a paragraph opens with a claim in bold and then four more
are appended with semicolons and dashes until the opening claim is the only one
anyone sees.

The rule is therefore that the claim is the first sentence, everything after it
supports that claim, and a paragraph making two claims is two paragraphs. The gate
lists every prose paragraph over 110 words or four sentences and fails while any
remain. Tables, fenced blocks, lists and headings are outside its scope, because
none of them is read as prose.

The gate has a second mode that matters more than the first. A structural pass over
a document can split a paragraph and quietly drop the clause it was splitting off,
and nothing about the result looks wrong. So the tool will compare two versions of
a file word for word and name anything that disappeared. That check caught five
genuinely lost words during the pass that introduced it.

\section{Why the Format Is Part of the System}

It would be possible to treat all of this as documentation hygiene. It is not, and
the reason returns to the paper's argument.

The studio's knowledge reaches an agent through exactly four channels. The
\textbf{API} tells it what may be asked. The \textbf{refusals} tell it what may not
be built, and \cref{ch:refusals} showed how much design is carried in them. The
\textbf{worked examples} tell it what is worth building. The \textbf{instructions}
tell it how to find out what it has actually built, which is the one thing an agent
authoring blind cannot discover by looking.

Only the last two have any opinion about quality, and neither of them runs.

Nothing in the first two channels would stop an agent building a board that used
one instrument and called it terrain. No rule fires. Every gate passes. The board
exports. The only thing that prevents it is having read an example that put five
treatments side by side and made the difference visible, and the only thing that
prevents the subtler version of the same failure is a card that says, in a section
of its own, which mistake this instrument invites.
